% !TeX program = lualatex
% !TeX encoding = utf8
% !TeX spellcheck = uk_UA
% !TeX root = ../main.tex

%=========================================================
\chapter{Електричне поле у провідниках}\label{\currfilebase}\hypertarget{FieldInMetalls}{}
\graphicspath{{\currfilebase/Pictures}}
%=========================================================
\epigraph{\Alexander Я зайшов у куб і жив у ньому, і, користуючись запаленими свічками, електрометрами та всіма
іншими засобами виявлення електричних станів, не зміг виявити жодного впливу на них\ldots\ хоча весь цей час
зовнішня поверхня куба була сильно заряджена, і великі іскри та китиці зривалися з усіх її частин.}{Майкл Фарадей,
\enquote{Експериментальні дослідження з електрики}, \S\,1174}


%% --------------------------------------------------------
\section{Структура провідників}\hypertarget{conductor-structure}{}
%% --------------------------------------------------------



Провідниками називають речовини, в структурі яких є вільні заряди (електрони в металах,
наприклад), які можуть переміщуватися під дією як завгодно слабких полів (рис.~\ref{tikz:metall_structure}).
\begin{wrapstuff}[o, top=0, type=figure, width=0.3\linewidth, belowsep=1em, abovesep=1em]
    \centering
    \localinput{metall_structure.tikz}
    \caption{Модель будови металу (електронний газ у кристалічній решітці)}
    \label{tikz:metall_structure}
\end{wrapstuff}


Електронні оболонки атомів, що складають кристалічну решітку металів, сильно
перекриваються, внаслідок чого не можна вказати, біля якого іона локалізовано той чи інший
\emphx{електрон валентної оболонки}, --- вони легко перетікають від одного іона до іншого,
і в цьому разі кажуть, що \emphx{електрони колективізовані}. Іони --- це ядра й
електрони внутрішніх оболонок, які сильно локалізовані. Електрони, які делокалізовані,
вільно переміщаються по кристалу. Саме вільні електрони відповідають за багато фізичних і,
особливо, транспортних властивостей металів.


%% --------------------------------------------------------
\subsection*{Досліди Толмена та Стюарта}
%% --------------------------------------------------------


У\index{Толмена і Стюарта досліди} 1916~р. американські фізики Р.~Толмен (1881--1948) і~Т.~Д.~Стюарт виконали
кількісні виміри, які неспростовно довели, що струм у металевих провідниках зумовлений рухом
вільних електронів (рис.~\ref{pic:TolmenStuartExperiment}).

%---------------------------------------------------------
\begin{SCfigure}[][h!]%[10]{O}{0.3\linewidth}
    \centering
    \includegraphics[width=0.33\linewidth]{exptommstuart}
    \caption{Принципова ідея дослідної установки Толмена і Стюарта}
    \label{pic:TolmenStuartExperiment}
\end{SCfigure}
%---------------------------------------------------------
У цих дослідах котушку з великим числом витків тонкого дроту підключали до
гальванометра і приводили в швидке обертання навколо своєї осі. Під час різкого
гальмування котушки в колі виникав короткочасний струм, зумовлений інерцією носіїв
заряду. За напрямком відхилення стрілки гальванометра було встановлено, що
\emph{електричний струм створюють негативно заряджені частинки}. При цьому
експериментально був отриманий питомий заряд носіїв $q/m$, близький до питомого
заряду електрона, отриманого з інших дослідів. Так було експериментально доведено,
що носіями вільних зарядів у металах є електрони.


%\href{https://www.youtube.com/watch?v=nVwquffMk44}{\color{blue}Відеодемотнстрація дослідів}

%% --------------------------------------------------------
\section{Провідник у полі}\hypertarget{conductor-in-field}{}
%% --------------------------------------------------------


Якщо провідник потрапляє в поле, електрони в ньому починають рухатися проти поля.
На одній частині поверхні провідника виступає негативний заряд --- ця поверхня стає
збагаченою електронами. На протилежній частині поверхні
електронів виявляється дещо менше, ніж потрібно для нейтралізації позитивного
іонного заряду кристалічної решітки, і ця частина поверхні виявляється зарядженою
позитивно. Позитивна і негативна частини поверхні створюють своє власне поле, за
напрямком протилежне зовнішньому. Обидва поля --- зовнішнє $\Efield_\text{ex}$ і
поле власних     поверхневих зарядів провідника ($\Efield_\text{in}$) точно
компенсують одне одного в усіх точках усередині провідника. Біля поверхні зовні лишається лише нормальна складова повного поля,
$E_n=4\pi\sigma$ (див. нижче). Схема перерозподілу зарядів у провіднику,
внесеному в зовнішнє електричне поле, показана на рис.~\ref{pic:MetallInField}.

Якби в якій-небудь ділянці провідника поля $\Efield_\text{ex}$ і $\Efield_\text{in}$
не компенсували одне одного, то в цій ділянці на електрони діяла б сила і йшов би
струм. Відсутність струму означає повну компенсацію полів.


%=========================================================
\begin{figure}[h!]\centering
    %---------------------------------------------------------
    \begin{subfigure}[t]{0.45\linewidth}\centering
        \localinput{ChargeRedistribution.tikz}
        \caption{Перерозподіл зарядів до встановлення рівноваги}
        % \label{fig:ChargeRedistribution}
    \end{subfigure}
    %---------------------------------------------------------
    \begin{subfigure}[t]{0.45\linewidth}\centering
        \localinput{ChargeRedistributionBefore.tikz}
        \caption{Розподіл заряду після встановлення рівноваги}
        % \label{fig:ChargeRedistributionBefore}
    \end{subfigure}
    %---------------------------------------------------------
    \caption{Провідник в електричному полі}
    \label{pic:MetallInField}
\end{figure}
%=========================================================


\emph{У стаціонарному стані}, коли немає струмів, \emph{в об'ємі
    провідника результуюче електричне поле}:
\begin{equation}
    \Efield = 0
\end{equation}
Струм буде текти доти, доки заряди не розташуються так, щоб поле було
відсутнє усюди в об'ємі провідника.


З теореми Гаусса
\begin{equation*}
    \Efield = 0\ {\color{red}\Rightarrow}\ \rho = \frac1{4\pi} \Div\Efield\
    {\color{red}\Rightarrow}\ \rho = 0,
\end{equation*}
об'ємна густина зарядів у провіднику дорівнює нулю, \emph{нескомпенсовані (надлишкові) заряди можуть
    розташовуватися тільки на поверхні}.


Граничні умови
\begin{align*}
    E_n      & = 4\pi\sigma, \\
    E_{\tau} & = 0
\end{align*}
показують, що \emph{силові лінії підходять до поверхні провідника перпендикулярно} і~закінчуються
(або починаються) на поверхневих зарядах: усередині провідника поля немає.


\begin{Attention}
    Поверхня провідника є еквіпотенціальною поверхнею! Потенціал в об'ємі провідника
    сталий і дорівнює потенціалу на його поверхні:
    \begin{equation*}
        \phi_S = \phi_V = \const.
    \end{equation*}
\end{Attention}


%% --------------------------------------------------------
\subsection*{Електрофорна машина (машина Вімсхерста)}
%% --------------------------------------------------------

Явище\index{Вімсхерста машина} перерозподілу вільних зарядів у провіднику під дією зовнішнього поля (розглянуте вище, у §\enquote{Провідник у полі}) лежить в основі \emph{електрофорної машини Вімсхерста} --- індукційного генератора, який, обертаючись, сам нарощує заряд до високих напруг завдяки позитивному зворотному зв'язку, без потреби заряджати щось ззовні тертям (рис.~\ref{pic:wimshurst}).

Машина складається з двох дисків із діелектрика (скла чи ебоніту), закріплених на спільній осі, які обертаються в протилежних напрямках; на зовнішньому боці кожного диска симетрично наклеєні радіальні металеві сектори. Біля кожного диска розташований \emph{нейтралізатор}~--- діагональний металевий стрижень зі щітками на обох кінцях, які торкаються двох діаметрально протилежних секторів цього диска; нейтралізатори двох дисків повернуті один відносно одного (зазвичай на кут близько $90^\circ$).

Заряд наростає завдяки індукції. Нехай сектор диска~А має невеликий випадковий додатний заряд. Коли він проходить навпроти сектора диска~Б, що саме торкається щітки нейтралізатора~Б, його поле притягує на цей сектор від'ємний заряд, а на з'єднаний із ним через нейтралізатор протилежний сектор~--- додатний. Обертаючись, заряджені сектори диска~Б проходять навпроти секторів диска~А біля нейтралізатора~А й так само індукують на них заряди, вже більші за величиною. Так з кожним обертом заряд наростає лавиноподібно (позитивний зворотний зв'язок), аж поки не встановиться рівновага між його приростом і втратами (іонізацією повітря, коронним розрядом тощо).

Накопичений заряд знімають гребінчасті збирачі біля країв дисків, розташовані там, де сектори проходять із найбільшим потенціалом, --- і подають на дві лейденські банки (конденсатори, видно збоку на рис.~\ref{pic:wimshurst}) та на вихідні електроди з кульками на кінцях довгих стрижнів. Коли різниця потенціалів між кульками перевищує напругу пробою повітряного проміжку, відбувається іскровий розряд.

\begin{SCfigure}[1][h!]\centering
    % TODO: на рисунку не видно дисків А та Б, про які йдеться в тексті; автор шукає кращий рисунок
    \includegraphics[width=0.5\linewidth]{electrophorus}
    \caption{Електрофорна машина Вімсхерста: два зустрічно обертові диски з металевими секторами, схрещені нейтралізуючі стрижні, лейденські банки-конденсатори та вихідні електроди для іскрового розряду.}
    \label{pic:wimshurst}
\end{SCfigure}


%% --------------------------------------------------------
\subsection*{Поле в порожнині провідника}
%% --------------------------------------------------------


Сумарне поле зовнішніх зарядів і зарядів на поверхні провідника в товщі провідника
дорівнює нулю. \emph{Чи будуть з'являтись заряди на внутрішній поверхні
    провідника, в якому є порожнина?}

%=========================================================
\begin{figure}[h!]\centering
    %---------------------------------------------------------
    \captionbox{Заряд у порожнині\label{tikz:charge_in_cavity}}[0.45\linewidth]{
        \centering
        \localinput{cavity1.tikz}
    }
    %---------------------------------------------------------
    \captionbox{Чи будуть заряди на внутрішній поверхні провідника?\label{tikz:charge_inside}}[0.45\linewidth]{
        \centering
        \localinput{cavity2.tikz}
    }
    %---------------------------------------------------------
\end{figure}
%=========================================================

Нехай у~\emph{порожнині незарядженого провідника є заряд} $q$
(рис.~\ref{tikz:charge_in_cavity}). Проведемо замкнену поверхню в~товщі металу навколо
порожнини. На ній $\Efield = 0$, тому за теоремою Гаусса повний заряд усередині неї
дорівнює нулю: на внутрішній поверхні оболонки індукується заряд $-q$. Оскільки оболонка
в~цілому нейтральна, на її зовнішній поверхні з'являється заряд $+q$. Розподіл цього
заряду по зовнішній поверхні визначається лише формою провідника й зовнішніми зарядами
і~не залежить від того, де саме в~порожнині розташований заряд $q$: поле цього заряду
й~індукованого заряду $-q$ поза внутрішньою поверхнею взаємно компенсується.

Якщо в~\emph{порожнині зарядженого провідника немає зарядів}, то чи будуть заряди на
внутрішній поверхні оболонки (рис.~\ref{tikz:charge_inside})? За теоремою Гаусса (та сама
поверхня в~товщі металу) їхній \emph{сумарний} заряд дорівнює нулю. Але, можливо, на одних
ділянках внутрішньої поверхні є додатний заряд, а на інших~--- від'ємний? Тоді в~порожнині
була б силова лінія, що йде від \enquote{$+$} до \enquote{$-$}. Замкнімо контур уздовж цієї
лінії через порожнину й назад через метал: на першій ділянці $\Efield\cdot d\vect{r} > 0$,
а в~металі $\Efield = 0$, тож $\oint_L \Efield\cdot d\vect{r} > 0$, що суперечить теоремі
про циркуляцію. Отже, зарядів на внутрішній поверхні немає, і~поля в~порожнині теж немає.

\begin{Attention}
    Поле в~порожнині замкненої провідної оболонки \emph{не залежить} від зарядів зовні
    оболонки: зовнішні поля в~порожнину не проникають. Навпаки, заряди в~порожнині не
    впливають на \emph{розподіл} заряду на зовнішній поверхні, але від них залежить
    \emph{повний} заряд, що там з'являється. Якщо оболонку заземлити, цей заряд стікає
    в~землю, і~внутрішня та зовнішня частини простору стають електрично повністю
    незалежними.
\end{Attention}



%% --------------------------------------------------------
\subsection*{Електростатичний захист}
%% --------------------------------------------------------

Електростатичний\index{Електростатичний захист} захист --- це явище, яке дає змогу екранувати електричне поле, створюючи зону безпеки всередині замкненої металевої оболонки. На
практиці для цього часто використовують металеві сітки: для статичних і~повільно змінних полів сітка з~комірками, малими порівняно з~розмірами
захищеної області, діє майже так само, як суцільна оболонка.

\begin{figure}[h!]%{0.49\linewidth}
    \centering
    \includegraphics[width=0.75\linewidth]{FaradayCage}
    \caption{Людина в клітці Фарадея повністю захищена від впливу зовнішніх електричних полів.\label{pic:man_in_faraday_cage}}
\end{figure}

Явище було відкрито Майклом Фарадеєм 1836 року. Він звернув увагу, що зовнішнє
електричне поле не може потрапити всередину металевої клітки. Електростатичний
захист потрібен там, де необхідно екранувати електроприлади від зовнішніх
електричних полів (наприклад, чутливі вимірювальні прилади, лабораторне обладнання,
екрановані кабелі). Ці прилади поміщають у~металевий корпус, який захищає
їх від зовнішніх електричних перешкод (рис.~\ref{pic:man_in_faraday_cage}). Для швидкозмінних
полів екранування теж працює, якщо товщина стінок більша за так звану глибину скін-шару;
про це йтиметься в~розділах про змінні поля.

Той самий ефект пояснює, чому \emph{автомобіль і літак є відносно безпечним місцем під час грози}: металевий кузов чи фюзеляж діє як клітка Фарадея, і при ударі блискавки струм тече по зовнішній провідній обшивці, не проникаючи всередину салону, де електричне поле залишається практично нульовим. З цієї ж причини кабіни літаків, корпуси суден і навіть металізовані намети для чутливого обладнання захищають вміст від зовнішніх електромагнітних полів і розрядів.

%\href{https://www.youtube.com/watch?v=jsBmv3-i2g0}{Відео: клітка Фарадея}




%% --------------------------------------------------------
\subsection*{Електростатичний генератор Ван де Граафа}
%% --------------------------------------------------------



Генератор\index{Ван де Граафа генератор} Ван де Граафа --- електростатичний генератор, в якому використовуються властивості
циліндра Фарадея для створення високої напруги. Таким методом можна досягнути напруги до
кількох мегавольт (рис.~\ref{pic:VandeGraaf}).

\begin{wrapstuff}[type=figure, width=0.5\linewidth, o, top=3, belowsep=1em]
    \centering
    \includegraphics[width=0.95\linewidth]{VandeGraaf}
    \caption{Генератор Ван де Граафа\label{pic:VandeGraaf}}
\end{wrapstuff}
Простий генератор Ван де Граафа складається з діелектричної (шовкової або гумової) стрічки ($4$ і $5$ на малюнку), що обертається на роликах $3$ і $6$,
причому верхній ролик діелектричний, а нижній металевий і з'єднаний із землею. Верхній ролик $3$ разом із верхньою частиною стрічки розміщений усередині порожнистої металевої сфери $1$. Два електроди $2$ і
$7$ у формі щіток розташовані на невеликій відстані від стрічки зверху і знизу, причому електрод $2$ з'єднаний з внутрішньою поверхнею сфери $1$. Поблизу
нижнього електрода повітря іонізується, утворюються позитивні іони, які під дією сили Кулона рухаються до заземленого ролика $6$ й осідають на стрічці,
завдяки чому
частина стрічки, що рухається вгору, заряджається. Стрічка доставляє заряд усередину сфери $1$, де його знімає щітка $2$. Тут і~полягає головна ідея:
заряд самої сфери розташований на її зовнішній поверхні і~всередині поля не створює, тому заряд переходить зі стрічки на сферу й одразу йде на
її зовнішню поверхню, \emph{хоч яким великим уже є потенціал сфери}. Таким чином на зовнішній поверхні сфери
накопичується дедалі більший заряд. Можливість отримання високої напруги обмежена коронним розрядом, що виникає при іонізації повітря навколо сфери.

\begin{Attention}
    Основний принцип роботи генератора базується на властивості заряду концентруватися на зовнішній поверхні провідника, тоді як усередині сфери поля
    немає. Це дає змогу передавати заряд зі стрічки на сферу, де він накопичується.
\end{Attention}

Сучасні генератори Ван де Граафа замість стрічок використовують ланцюги, що складаються з металевих і пластикових ланок, що чергуються між собою.



%\href{https://www.youtube.com/watch?v=Xqt2gAalV4Y}{Відео: Принцип дії генератора Ван де
%	Граафа}





%% --------------------------------------------------------
\section{Електричне поле біля вістря}\hypertarget{field-near-tip}{}
%% --------------------------------------------------------

%Електричні заряди завжди розподіляються тільки на поверхні провідника. Однак цей розподіл різний залежно від форми провідника.

Електричні\index{Вістря, поле поблизу} заряди завжди розподіляються тільки на поверхні провідника. Однак цей розподіл різний залежно від форми провідника: \emph{напруженість
    електричного поля і поверхнева густина заряду мають найбільші значення біля загострення провідника}. Покажемо це на найпростішій моделі.

\begin{figure}[h!]
    \centering
    \localinput{two_spheres_wire.tikz}
    \caption{Дві провідні кулі різних радіусів, з'єднані тонким дротом: обидві мають однаковий потенціал $\phi$, проте на меншій кулі густина заряду й поле біля поверхні більші.\label{tikz:two_spheres_wire}}
\end{figure}

Розглянемо дві провідні кулі різних радіусів $R_1$ і $R_2$ ($R_1 < R_2$), з'єднані тонким провідним дротом і розташовані далеко одна від одної --- настільки, що взаємним впливом їхніх полів одна на одну (а також зарядом самого дроту) можна знехтувати (рис.~\ref{tikz:two_spheres_wire}). Тоді заряд кожної кулі можна вважати розподіленим по її поверхні рівномірно, як у відокремленої кулі.

Оскільки кулі з'єднані провідником, увесь цей провідник (обидві кулі й дріт) є єдиним провідником у стані рівноваги, а отже, еквіпотенціальним:
\begin{equation*}
    \phi_1 = \phi_2 = \phi.
\end{equation*}
Потенціал відокремленої зарядженої кулі радіуса $R$ із зарядом $Q$ дорівнює $\phi = Q/R$. Виразимо звідси заряд через потенціал і радіус, а тоді знайдемо поверхневу густину заряду:
\begin{equation*}
    Q = \phi R, \qquad \sigma = \frac{Q}{4\pi R^2} = \frac{\phi}{4\pi R}.
\end{equation*}
Отже, для обох куль
\begin{equation*}
    \sigma_1 R_1 = \sigma_2 R_2 = \frac{\phi}{4\pi} = \const,
\end{equation*}
тобто поверхнева густина заряду обернено пропорційна радіусу кулі: $\sigma \sim 1/R$. Оскільки напруженість поля біля поверхні провідника $E = 4\pi\sigma$, вона також обернено пропорційна радіусу:
\begin{equation*}
    \frac{E_1}{E_2} = \frac{\sigma_1}{\sigma_2} = \frac{R_2}{R_1} > 1,
\end{equation*}
тобто на кулі меншого радіуса і густина заряду, і поле сильніші.

\begin{Attention}
    Результат отримано для двох далеких куль (їхньою взаємодією знехтувано), тому для провідника довільної форми він дає лише \emph{якісне} правило: густина заряду зростає зі збільшенням кривизни поверхні, тобто обернено залежить від \emph{локального радіуса кривизни}, але точної пропорційності $\sigma\sim1/R$ у кожній точці немає. Саме тому там, де провідник загострюється (радіус кривизни малий), густина заряду і поле різко зростають, а на плоских чи слабо викривлених ділянках (великий радіус кривизни) --- значно менші (рис.~\ref{pic:tip}).
\end{Attention}


\begin{SCfigure}[1][h!]%{0.39\linewidth}
    \centering
    \includegraphics[width=0.4\linewidth]{tip.png}
    \caption{Поле біля вістря}
    \label{pic:tip}
\end{SCfigure}

%% --------------------------------------------------------
\subsection*{Іонний вітер}
%% --------------------------------------------------------

\enquote{Іонний вітер}\index{Іонний вітер} --- явище, за якого рух повітря створюється за допомогою
електричного поля.

Наявні в повітрі в невеликій кількості вільні заряди (іони обох знаків і електрони)
поблизу вістря розганяються сильним полем і, вдаряючись об молекули газу, іонізують
їх. Створюється область просторового заряду, звідки іони того ж знака, що й вістря,
виштовхуються полем, захоплюючи за собою молекули газу. Потік молекул та іонів,
спрямований від вістря, створює враження \enquote{стікання зарядів}.
При цьому вістря розряджається (іонами протилежного знака, що потрапляють на
нього) і одночасно отримує імпульс, спрямований проти стікання (рис.~\ref{pic:IonWind}).

\begin{figure}[h!]
    \centering
    \includegraphics[width=0.95\linewidth]{ion_wind}
    \caption{Ілюстрація іонного вітру}
    \label{pic:IonWind}
\end{figure}

Отже, загострений провідник у сильному полі не просто \enquote{притягує}
розряд --- він активно \emph{змінює} повітря навколо себе: іонізує його
й жене іони разом із молекулами газу від вістря. Це явище саме собою
цікаве, але воно має й цілком практичне застосування: якщо підняти
загострений заземлений металевий стрижень над будівлею, ми отримаємо
пристрій, відомий під назвою \emph{громовідвід}.

%\href{https://www.youtube.com/watch?v=LNobNwDX2jQ}{Відео: Демонстрація іонного вітру}

%% --------------------------------------------------------
\subsection*{Громовідвід (блискавковідвід)}
%% --------------------------------------------------------

Концентрація\index{Громовідвід} поля біля загострень і пов'язаний з нею іонний вітер
(див. попередній підрозділ) лежать в основі \emph{громовідводу},
винайденого Б.~Франкліном 1752 року. Пристрій складається з
металевого стрижня, піднятого над найвищою точкою будівлі, товстого
провідника (\emph{струмовідводу}), що з'єднує стрижень із землею
вздовж стіни споруди, і заземлювача --- металевого електрода,
закопаного в ґрунт (рис.~\ref{pic:lightning_rod}).

\begin{SCfigure}[1][h!]\centering
    \includegraphics[width=0.5\linewidth]{lightning_rod}
    \caption{Громовідвід: стрижень на даху будівлі, струмовідвід уздовж стіни та заземлювач.}
    \label{pic:lightning_rod}
\end{SCfigure}

Громовідвід захищає будівлю двома способами. Перший ---
\emph{послаблення заряду хмари заздалегідь}. Біля вістря стрижня
поле грозової хмари стає набагато сильнішим, ніж навколо, і повітря
починає іонізуватися --- виникає той самий іонний вітер, який ми
щойно розглянули. Частина заряду хмари над будівлею поступово
нейтралізується, і поле не встигає (або майже не встигає) досягти
тієї величини, за якої повітря пробивається і виникає блискавка над
самою спорудою.

Другий --- \emph{перехоплення розряду}. Якщо блискавка все ж починає
розвиватися, то найімовірніше вона \enquote{вибере} саме вістря
стрижня, бо там поле найсильніше. Від вістря назустріч розряду
піднімається зустрічний канал іонізованого повітря; з'єднавшись,
вони утворюють шлях блискавки, який веде до громовідводу, а не до
будівлі. Далі струм по струмовідводу й заземлювачу йде прямо в землю,
і конструкція та все, що в ній, не нагріваються й не займаються.

\begin{Historical}[unbreakable]
    Історично Франклін вважав головним саме перший механізм ---
    поступове \enquote{розряджання} хмари. Сучасні дослідження
    показують, що на практиці основний захист дає саме
    \emph{перехоплення розряду}: внесок превентивного ефекту
    в реальне зниження ймовірності удару блискавки значно менший,
    ніж думали раніше, і це досі є предметом наукової дискусії.
\end{Historical}

Область простору навколо громовідводу, куди практично не влучає блискавка (оскільки
розряд \enquote{обирає} ближчий і сильніший осередок поля --- сам стрижень), називають
\emph{зоною захисту}; її розміри й форма залежать від висоти стрижня і використовуються
під час проєктування блискавкозахисту будівель.


% --------------------------------------------------------
\section{Єдиність розв'язку електростатичної задачі}\hypertarget{uniqueness-theorem}{}
%% --------------------------------------------------------

Задача електростатики зводиться до \emph{знаходження розв'язку}
рівняння Пуассона~\eqref{eq:Puasson} (або Лапласа
\eqref{eq:Laplase}). Нагадаємо (\S~\ref{sec:Uniqueness}), що за \emph{теоремою єдиності}
розв'язок, який задовольняє рівняння й~граничні умови, єдиний. Тому розв'язок задачі можна
\emph{вгадати}: якщо вгадана функція в~деякій області задовольняє рівняння Лапласа (або
Пуассона) і~\emph{умови на границі області}, то вона і~є правильним розв'язком.

Для задач із провідниками граничні умови мають простий вигляд: на поверхні кожного
провідника потенціал сталий, і~задано \emph{або} сам цей потенціал (зокрема, $\phi = 0$ для
заземленого провідника), \emph{або} повний заряд ізольованого провідника. В~обох випадках
розв'язок єдиний.

\begin{Example}[unbreakable]{Доведення через тотожність Гріна}
    Це доведення доповнює наведене в~\S~\ref{sec:Uniqueness} і~придатне як для заданого на межі потенціалу, так і~для заданої нормальної складової поля.

    Нехай $\phi_1$ і $\phi_2$ --- два розв'язки рівняння Пуассона з однією й тією ж густиною заряду $\rho$ в об'ємі $V$, що задовольняють однакові умови на межі $S$ цього об'єму. Розглянемо їхню різницю $\psi = \phi_1 - \phi_2$. Вона задовольняє рівняння Лапласа $\nabla^2\psi = 0$ в $V$, а на $S$ або $\psi = 0$ (задано потенціал), або $\partial\psi/\partial n = 0$ (задано $E_n$).

    Скориставшись першою тотожністю Гріна~\eqref{Grin1} (теоремою про дивергенцію, застосованою до $\psi\Grad\psi$):
    \begin{equation*}
        \int\limits_V (\Grad\psi)^2\, dV = \oint\limits_S \psi\,\Grad\psi \cdot d\vect{S} - \int\limits_V \psi\, \nabla^2\psi\, dV,
    \end{equation*}
    бачимо, що обидва доданки праворуч зникають: другий --- бо $\nabla^2\psi = 0$ в $V$, перший --- бо на $S$ дорівнює нулю $\psi$ або $\Grad\psi\cdot d\vect{S} = (\partial\psi/\partial n)\,dS$. Отже,
    \begin{equation*}
        \int\limits_V (\Grad\psi)^2\, dV = 0.
    \end{equation*}
    Підінтегральна функція невід'ємна, тому рівність нулю інтеграла можлива лише за умови $\Grad\psi = 0$ в усьому об'ємі $V$, тобто $\psi = \const$: поля $\Efield_1$ і~$\Efield_2$ збігаються. Якщо на межі задано потенціал, то $\psi = 0$ на $S$, і~ця стала дорівнює нулю: $\phi_1 \equiv \phi_2$. Розв'язок єдиний.
\end{Example}

Використання теореми єдиності дуже \emph{спрощує} розв'язання деяких електростатичних
задач.



%% --------------------------------------------------------
\subsection*{Метод електричних зображень}
%% --------------------------------------------------------


Метод\index{Електричних зображень метод} \emphx{ґрунтується на теоремі єдиності} і \emph{полягає у вгадуванні конфігурації зарядів}, у якій поле в цікавій для нас частині простору було б
тим же самим, що й у розглядуваній задачі.


Розглянемо ідею цього методу на прикладі, коли точковий заряд $q$
знаходиться біля нескінченної заземленої провідної площини (рис.~\ref{tikz:charge_near_plane}).

\begin{SCfigure}[][h!]
    \centering
    \includegraphics[width=6cm]{charge_near_plane.pdf}
    \caption{Точковий заряд $q$ поблизу нескінченної заземленої провідної площини.\label{tikz:charge_near_plane}}
\end{SCfigure}


Як знайти поле у верхній частині простору? Як знайти розподіл заряду на площині?

Можна вгадати таку конфігурацію зарядів, у якої є еквіпотенціальна поверхня з~$\phi = 0$,
що збігається з~площиною.

Для обчислення поля достатньо ввести \emph{фіктивний заряд-зоб\-ра\-жен\-ня $q' = -q$},
помістивши його по інший бік площини на такій самій відстані від неї, що й заряд $q$.
Дійсно, для будь-якої точки на самій площині відстані до заряду $q$ і до заряду-зображення
$q'$ однакові (позначимо цю спільну відстань $r$) --- через дзеркальну симетрію їхнього
розташування відносно площини. Тому внески зарядів у потенціал скорочуються:
\begin{equation*}
    \phi = \frac{q}{r} + \frac{q'}{r} = \frac{q}{r} - \frac{q}{r} = 0,
\end{equation*}
і площина справді виявляється еквіпотенціальною поверхнею з $\phi = 0$ --- саме такою, якою
має бути заземлена провідна площина.

Фіктивний заряд $q'$ створює у верхньому півпросторі точно таке саме поле, як і індуковані
заряди на площині (див. рис.~\ref{pic:image_method}). Зокрема, якщо заряд $q$ перебуває на
відстані $d$ від площини, то він притягується до неї з~силою, що дорівнює силі притягання
до заряду-зображення: $F = q^2/(2d)^2$.

\begin{Attention}
    Заряд-зображення відтворює правильне поле \emph{лише в тій половині простору, де
        знаходиться реальний заряд $q$} (тобто там, де немає провідника). У другій половині
    --- усередині провідника --- реальне поле, як показано раніше, дорівнює нулю; поле,
    яке ніби створює там заряд-зображення, --- суто математична фікція, потрібна лише для
    підбору правильного розв'язку в фізичній області, і жодного стосунку до дійсного поля
    всередині провідника не має.
\end{Attention}

Ідея заряду-зображення лежить в основі опису взаємодії зонда з провідною поверхнею в \emph{електростатичній силовій мікроскопії}: заряджений металевий зонд-вістря, піднесений близько до провідного зразка, індукує на поверхні заряд-зображення, і~сила його притягання до поверхні (яку можна оцінити саме методом зображень) використовується для побудови зображення поверхні з нанометровою роздільною здатністю.


%---------------------------------------------------------
\begin{figure}[h!]\centering
    \begin{subfigure}{0.45\linewidth}
        \centering
        \includegraphics[width=6cm]{two_charges.pdf}
        \caption{Конфігурація поля двох різнойменних зарядів. }
    \end{subfigure}
    \begin{subfigure}{0.45\linewidth}
        \centering
        \includegraphics[width=6cm]{charge_image.pdf}
        \caption{Заряд над площиною та його зображення.}%
    \end{subfigure}
    \caption{Над площиною картини поля на обох рисунках однакові.}
    \label{pic:image_method}
\end{figure}
%---------------------------------------------------------


%% --------------------------------------------------------
\subsection*{Заряд поблизу сферичної поверхні}
%% --------------------------------------------------------


Розглянемо систему двох різнойменних точкових зарядів, \emph{різних за модулем}. Можна
показати, що її еквіпотенціальна поверхня $\phi = 0$ є сферою (для рівних за модулем
зарядів вона вироджується в~площину). Це дає
змогу застосувати метод зображень до задачі про заряд поблизу провідної сфери (рис.~\ref{pic:mirror_sphere}).


\begin{figure}[h!]
    \centering
    \begin{subfigure}{0.45\linewidth}
        \centering
        \includegraphics[width=\linewidth]{mirrorsphere_fieldlines.pdf}
        \caption{Конфігурація силових ліній.}
        \label{pic:mirrorsphere2}
    \end{subfigure}\hfill
    \begin{subfigure}{0.45\linewidth}
        \centering
        \includegraphics[width=\linewidth]{mirrorsphere_scheme.pdf}
        \caption{До розрахунку положення і величини заряду-зображення.}
        \label{pic:mirrorsphere1}
    \end{subfigure}
    \caption{Заряд поблизу заземленої сфери: \protect\subref{pic:mirrorsphere2}~---
    конфігурація силових ліній; \protect\subref{pic:mirrorsphere1}~--- геометрія
    для розрахунку положення і величини заряду-зображення.}
    \label{pic:mirror_sphere}
\end{figure}

Нехай точковий заряд $q$ розташований на відстані $b$ від центра \emph{заземленої}
провідної сфери радіуса $R$ (рис.~\ref{pic:mirrorsphere2}). Поле поза сферою можна знайти, якщо підібрати такий
фіктивний заряд-зображення $q'$, щоб поверхня сфери була еквіпотенціальною з
$\phi = 0$. За теоремою єдиності знайдений так розв'язок і~є правильним.

З~міркувань симетрії заряд $q'$ лежить на прямій, що з'єднує центр сфери із зарядом $q$,
на відстані $a$ від центра. Вимагаючи $\phi = 0$ у~найближчій і~найвіддаленішій від $q$
точках сфери, маємо
\begin{equation*}
    \frac{q}{b - R} + \frac{q'}{R - a} = 0, \qquad \frac{q}{b + R} + \frac{q'}{R + a} = 0.
\end{equation*}
Розв'язуючи цю систему, отримаємо
\begin{equation}\label{eq:mirror_sphere}
    a = \frac{R^2}{b}, \qquad q' = -\frac{R}{b}\,q .
\end{equation}
Можна перевірити, що за цих значень $\phi = 0$ не лише у~двох вибраних точках, а~й~на всій
сфері. Повний індукований на сфері заряд дорівнює $q'$ (за теоремою Гаусса), а~сила
притягання заряду $q$ до сфери дорівнює силі його взаємодії із зарядом-зображенням:
$F = qq'/(b - a)^2$.
Заряд-зображення лежить на тій самій прямій, що й реальний заряд $q$, але з
протилежного боку від центра сфери, всередині неї. Знак мінус означає, що
індукований на сфері заряд має протилежний знак до $q$.

Якщо сфера \emph{ізольована} й \emph{несе повний заряд $q_0$}, то одного
заряду-зображення вже недостатньо: потенціал ізольованої сфери вже не дорівнює нулю, а~її
повний заряд має дорівнювати $q_0$. Тому до заряду $q'$ треба додати другий
заряд-зображення $q''$, розміщений у центрі сфери. Заряд, поміщений точно в~центр
сфери, створює на її поверхні сталий потенціал $q''/R$, тож поверхня залишається
еквіпотенціальною. Величина $q''$ визначається з умови, що повний заряд сфери (а~за
теоремою Гаусса він дорівнює сумі зарядів-зображень усередині неї) дорівнює $q_0$:
\begin{equation*}
    q' + q'' = q_0 .
\end{equation*}
Отже, поле поза ізольованою зарядженою сферою еквівалентне полю двох фіктивних
зарядів: $q'$ на відстані $a$ від центра (рис.~\ref{pic:mirrorsphere1}) і $q''$ у
центрі сфери. Поле всередині провідної сфери при цьому дорівнює нулю, а поле
заряду-зображення $q'$ у тій області є суто математичною фікцією, яка не має
фізичного змісту.


%%% --------------------------------------------------------
%\section{Металева куля в однорідному полі}
%%% --------------------------------------------------------
%
%Металева куля набуває дипольного моменту внаслідок перерозподілу зарядів по її поверхні:
%\begin{equation*}
%	\vect{p} = R^3\Efield_0
%\end{equation*}
%
%
%\begin{figure}[h!]
%\centering
%\localinput{EfieldMetall.tikz}
%\caption{Конфігурація поля навколо металевої кулі, вміщеної в однорідне електричне поле.\label{tikz:EfieldMetall}}
%\end{figure}
%
%
%
%
%Потенціал кулі:
%\begin{equation*}
%	\phi =
%	\begin{cases}
%		0,                                                                      & r \le R \\
%		-\left( \Efield_0\cdot \vect{r}\right) + \frac{\vect{p} \vect{r}}{r^3}, & r > R
%	\end{cases}
%\end{equation*}
%
%
%
%Поле зовні та всередині кулі:
%\begin{equation*}
%    \Efield =
%    \begin{cases}
%    0,
%    & r \le R \\
%    \Efield_0 - \frac{\vect{p}}{r^3} + \frac{3\left(\vect{p}\vect{r}\right)\vect{r}
%    }{r^5}, & r > R,
%    \end{cases}
%\end{equation*}
%
%
%Заряди на поверхні:
%\begin{equation*}
%   			\sigma = \frac{3}{4\pi} \frac{\Efield_0\vect{r}}{R}.
%\end{equation*}

%% --------------------------------------------------------
\section{Металева куля в однорідному полі}\hypertarget{metal-sphere-in-field}{}
%% --------------------------------------------------------

\subsection*{Постановка задачі}

Металеву кулю радіуса $R$ вміщено в однорідне зовнішнє
електричне поле $\Efield_0$. Метал --- провідник, тому в стані
рівноваги вільні заряди в ньому нерухомі, а отже, поле всередині
дорівнює нулю:
\begin{equation*}
    \Efield = 0, \qquad r < R .
\end{equation*}
Індуковані заряди розташовуються на поверхні кулі (у~цілому куля нейтральна). Потрібно знайти
поле зовні, розподіл заряду на поверхні та наведений дипольний момент.

Задача має аксіальну симетрію навколо напрямку $\Efield_0$, тому
природно вибрати сферичні координати з полярною віссю вздовж
$\Efield_0$, а початок координат сумістити з центром кулі (рис.~\ref{fig:sphere_in_field}).

\begin{figure}[h!]
    \centering
    \localinput{sphere_in_field.tikz}
    \caption{Металева куля радіуса $R$ в однорідному зовнішньому полі
        $\Efield_0$. Початок координат $O$ суміщено з центром кулі, полярну
        вісь напрямлено вздовж $\Efield_0$. Положення точки $P$ задається
        радіус-вектором $\vect{r}$ та кутом $\theta$ між $\vect{r}$ і
        полярною віссю.}
    \label{fig:sphere_in_field}
\end{figure}

%%% --------------------------------------------------------
\subsection*{Граничні умови на поле}
%%% --------------------------------------------------------


На поверхні металу $r = R$ мають виконуватися дві умови для поля.

\emph{Перша: тангенціальна складова неперервна.} Це загальна умова
для будь-якої межі розділу, і вона випливає з $\oint \Efield\cdot d\vect{r} = 0$:
\begin{equation*}
    E_{\tau}(R+0) = E_{\tau}(R-0) .
\end{equation*}
Оскільки всередині поля немає, $E_{\tau}(R-0) = 0$, а отже,
\begin{equation*}
    E_{\tau}(R+0) = 0 .
\end{equation*}
Це означає, що зовні поле підходить до поверхні кулі строго
перпендикулярно до неї.

\emph{Друга: стрибок нормальної складової.} Він визначається
поверхневим зарядом і випливає з теореми Гаусса:
\begin{equation*}
    E_{r}(R+0) - E_{r}(R-0) = 4\pi\sigma.
\end{equation*}
Оскільки $E_r(R-0) = 0$, маємо
\begin{equation*}
    \sigma = \frac{1}{4\pi}\,E_r(R+0) .
\end{equation*}

Зауважимо, що друга умова не обмежує поле: густина $\sigma$ заздалегідь
невідома і~просто визначається вже знайденим полем. Отже, задача
зводиться до такого: знайти поле зовні кулі, яке на її поверхні не
має дотичної складової ($E_\theta(R) = 0$), а~на великих відстанях
переходить в~однорідне поле $\Efield_0$. Коли таке поле знайдено,
друга умова дає розподіл заряду.

%%% --------------------------------------------------------
\subsection*{Поле зовні кулі}
%%% --------------------------------------------------------


Зовні кулі поле складається із зовнішнього однорідного поля
$\Efield_0$ та поля наведених зарядів на поверхні. Наведений заряд
в цілому нейтральний, тому його поле на великих відстанях
зменшується як поле диполя. Тому спробуємо (\enquote{вгадаємо}) розв'язок у~вигляді
\begin{equation*}
    \Efield(r,\theta) = \Efield_0 + \Efield_{\text{дип}},
\end{equation*}
де $\Efield_{\text{дип}}$ --- поле точкового диполя, розташованого в~центрі кулі й
орієнтованого вздовж $\Efield_0$, з невідомим поки що моментом $\vect{p}$. Якщо вдасться
підібрати $\vect{p}$ так, щоб виконувалися граничні умови, то за теоремою єдиності цей
розв'язок буде \emph{точним} у~всьому просторі поза кулею, а~не лише на великих
відстанях.

У сферичних координатах зручно працювати з компонентами. Для
однорідного поля, напрямленого вздовж осі $z$,
\begin{equation*}
    E_{0,r} = E_0\cos\theta, \qquad
    E_{0,\theta} = -E_0\sin\theta .
\end{equation*}
Для диполя з моментом $\vect{p} = p\,\vect{e}_z$ (задача~\ref{prb:DipolePolar}):
\begin{equation*}
    E_{\text{дип},r} = \frac{2p\cos\theta}{r^3}, \qquad
    E_{\text{дип},\theta} = \frac{p\sin\theta}{r^3} .
\end{equation*}

Повне поле зовні:
\begin{equation*}
    E_r = E_0\cos\theta + \frac{2p\cos\theta}{r^3}, \qquad
    E_\theta = -E_0\sin\theta + \frac{p\sin\theta}{r^3} .
\end{equation*}

%%% --------------------------------------------------------
\subsection*{Застосування граничних умов}
%%% --------------------------------------------------------


Підставляємо $r = R$.

З умови $E_\theta(R) = 0$:
\begin{equation*}
    -E_0\sin\theta + \frac{p\sin\theta}{R^3} = 0
    \quad\Longrightarrow\quad
    p = E_0 R^3 .
\end{equation*}

З умови $\sigma = E_r(R)/4\pi$:
\begin{equation*}
    \sigma(\theta) = \frac{1}{4\pi}
    \left( E_0\cos\theta + \frac{2p\cos\theta}{R^3} \right)
    = \frac{1}{4\pi} \left( E_0\cos\theta
    + 2E_0\cos\theta \right)
    = \frac{3}{4\pi}\,E_0\cos\theta .
\end{equation*}

Отже, умова $E_\theta(R) = 0$ визначила наведений дипольний момент, а~поле на поверхні~---
розподіл поверхневого заряду. Поле біля поверхні $E_r(R) = 3E_0\cos\theta$ найбільше
на полюсах кулі ($\theta = 0,\,\pi$), де воно втричі перевищує зовнішнє: саме там у~сильних
полях насамперед починається пробій повітря.

%%% --------------------------------------------------------
\subsection*{Остаточні результати}
%%% --------------------------------------------------------


Дипольний момент кулі:
\begin{equation*}
    \vect{p} = R^3\Efield_0 .
\end{equation*}

Поле зовні та всередині кулі:
\begin{equation*}
    \Efield =
    \begin{cases}
        0, & r \le R, \\[6pt]
        \Efield_0 - \dfrac{\vect{p}}{r^3}
        + \dfrac{3\left(\vect{p}\cdot\vect{r}\right)\vect{r}}{r^5},
           & r > R .
    \end{cases}
\end{equation*}

Заряди на поверхні:
\begin{equation*}
    \sigma = \frac{3}{4\pi}\,E_0\cos\theta
    = \frac{3}{4\pi}\,\frac{\Efield_0\cdot\vect{r}}{R}
    \quad\text{(у гауссовій системі)} .
\end{equation*}

Потенціал поля (для повноти; його легко відновити інтегруванням поля, а~сталу вибрано так, щоб потенціал кулі дорівнював нулю):
\begin{equation*}
    \phi =
    \begin{cases}
        0,                                    & r \le R, \\[4pt]
        -\left(\Efield_0\cdot\vect{r}\right)
        + \dfrac{\vect{p}\cdot\vect{r}}{r^3}, & r > R .
    \end{cases}
\end{equation*}

Отже, всередині металевої кулі поле відсутнє, а зовні воно є суперпозицією
зовнішнього однорідного поля $\Efield_0$ і поля індукованого диполя $\vect{p}$
(рис.~\ref{tikz:EfieldMetall}).

\begin{figure}[h!]
    \centering
    \localinput{EfieldMetall.tikz}
    \caption{Конфігурація поля навколо металевої кулі, вміщеної
        в однорідне електричне поле.}%
    \label{tikz:EfieldMetall}
\end{figure}

%% --------------------------------------------------------
\section{Сили, що діють на поверхню провідника в полі}\hypertarget{conductor-surface-force}{}
%% --------------------------------------------------------


На\index{Провідник!сили, що діють на поверхню} малий елемент $dS$ поверхні провідника діє сила:
\begin{equation*}
    dF = dq\, E' = \sigma dS \cdot  E',
\end{equation*}
з боку \emph{поля $E'$ інших зарядів} у місці, де
перебуває заряд $dq = \sigma dS$ (власне поле площадки на неї саму не діє). Але прилад
може виміряти лише \emph{сумарне поле $E$}, тобто поле \emph{інших зарядів $E'$} плюс
\emph{власне поле площадки $E_{\sigma}$}, і~розрізнити їх не може!

\begin{SCfigure}[1.5][h!]%{0.34\linewidth}
    \centering
    \localinput{Fmetall.tikz}
    \caption{До виведення сили, що діє на поверхню провідника}
\end{SCfigure}

Задача полягає в тому, щоб виразити силу саме через сумарне поле.

Мала площадка поблизу себе діє як нескінченна заряджена площина: її власне поле дорівнює
$E_\sigma = 2\pi\sigma$ і~напрямлене від площадки в~обидва боки (для $\sigma > 0$)~---
назовні з~одного боку і~вглиб провідника з~іншого. Поле $E'$ інших зарядів у~цьому місці
однакове з~обох боків площадки. Усередині провідника сумарне поле дорівнює нулю:
\begin{equation*}
    E' - E_{\sigma} = 0 \quad\Longrightarrow\quad E' = E_{\sigma} = 2\pi\sigma,
\end{equation*}
тобто \emph{густина заряду підлаштовується під поле інших зарядів}! Зовні поля
додаються, і~гранична умова дає поле у~вакуумі біля поверхні:
\begin{equation*}
    E = E' + E_{\sigma} = 2 E' = 4\pi\sigma.
\end{equation*}

Отже, сила, що діє на одиницю поверхні провідника, дорівнює
\begin{equation*}
    \frac{dF}{dS} = f = \sigma E' = \frac12 \sigma E = \frac{E^2}{8\pi}.
\end{equation*}
Вона залежить лише від сумарного поля біля поверхні й~завжди напрямлена \emph{назовні}
від провідника, незалежно від знака $\sigma$: поле ніби \enquote{розтягує} поверхню
провідника.


%% --------------------------------------------------------
\section*{Підсумки}
\phantomsection
\addcontentsline{toc}{section}{Підсумки}
%% --------------------------------------------------------

Провідник --- це речовина з вільними зарядами; в металах ними є \emphx{колективізовані} валентні електрони, що підтверджують досліди Толмена та Стюарта. У стані рівноваги поле в об'ємі провідника дорівнює нулю, $\Efield = 0$, тому об'ємна густина заряду також нульова: надлишковий заряд розташовується лише на поверхні. Поверхня провідника є \emphx{еквіпотенціальною}, а на ній виконуються граничні умови $E_n = 4\pi\sigma$, $E_\tau = 0$, тобто силові лінії підходять до неї перпендикулярно. Замкнена провідна оболонка електрично розділяє внутрішній і зовнішній простір: у порожнині без зарядів поля немає (це випливає з теореми про циркуляцію), а заряд усередині порожнини наводить на зовнішній поверхні такий самий заряд. На цьому ґрунтуються електростатичний захист (клітка Фарадея) і генератор Ван де Граафа, а перерозподіл зарядів у зовнішньому полі лежить в основі машини Вімсхерста.

Густина заряду й поле біля поверхні зростають зі збільшенням кривизни ($\sigma \sim 1/R$, $E = 4\pi\sigma$), звідки іонний вітер біля вістря й дія громовідводу. Розв'язок рівнянь Пуассона й Лапласа з заданими граничними умовами \emph{єдиний}; цей факт дає змогу розв'язок \emph{вгадувати}, що й робить метод електричних зображень. Для заземленої площини заряд-зображення $q' = -q$ стоїть дзеркально, а для заземленої сфери $a = R^2/b$, $q' = -Rq/b$; для ізольованої зарядженої сфери додається другий заряд-зображення $q''$ у центрі, причому $q' + q'' = q_0$. Металева куля в однорідному полі набуває дипольного моменту $\vect{p} = R^3\Efield_0$; поле всередині неї нульове, а поверхневий заряд $\sigma = \frac{3}{4\pi}E_0\cos\theta$. На поверхню провідника діє напрямлена назовні сила з густиною $f = \frac12\sigma E = E^2/8\pi$, яка залежить лише від сумарного поля.

У провіднику заряди рухаються вільно й повністю екранують внутрішнє поле. Більшість речовин, однак, --- діелектрики, у яких заряди зв'язані й можуть лише незначно зміщуватися під дією поля. Як таке зміщення (\emphx{поляризація}) змінює поле всередині речовини, розглянемо в наступному розділі (розділ~\ref{FieldInDielectrics}).

%% --------------------------------------------------------
\section*{Задачі}
\phantomsection
\addcontentsline{toc}{section}{Задачі}
%% --------------------------------------------------------

% задачі цього розділу короткі: не розриваємо їх між сторінками
\begingroup
\tcbset{answersolutionstyle/.append style={unbreakable}}

%=========================================================
\begin{problem}%
    Точковий заряд $q$ розташований на відстані $d$ від нескінченної заземленої металевої
    площини. Знайдіть густину $\sigma(x)$ індукованого заряду на площині, де $x$~---
    відстань уздовж площини від точки, що лежить навпроти заряду. Чому дорівнює повний
    індукований заряд?
\end{problem}

%=========================================================
\begin{problem}%
    Знайдіть потенціал \emph{незарядженої} і \emph{незаземленої} металевої сфери радіуса $R$,
    що перебуває в полі заряду $q$, розташованого на відстані $b$ ($b > R$) від центра сфери.
\end{problem}

%=========================================================
\begin{problem}%
    На відстані $b= 10R$ від заземленої незарядженої металевої сфери радіусом $R$ розташований
    точковий електричний диполь з моментом $p$, причому вісь диполя перпендикулярна прямій, що
    сполучає центр сфери з серединою осі диполя (див. рисунок). Знайдіть силу взаємодії між
    диполем і сферою.

    %Відповідь: $F = \frac{3p^2R^3}{b^7}$, взаємодія~--- притягування.

    \begin{center}
        \begin{tikzpicture}
            \draw [ball color=gray, opacity=0.5] (0,0) circle (1);
            \draw [-latex] (4,-0.5) -- (4,0.5) node [above] {$\vect{p}$};
            \draw [-latex'] (0,0) -- node [right] {$R$}(135:1);
            \draw [latex'-latex'] (0,0) -- node [below] {$b$}(4,0);
            \draw (0,-1) -- (0,-1.5) coordinate (G);
            \draw (G) +(-0.5,0) -- +(0.5,0)
            +(-0.3,-0.1) -- +(0.3,-0.1)
            +(-0.1,-0.2) -- +(0.1,-0.2);
        \end{tikzpicture}
    \end{center}
\end{problem}

%=========================================================
\begin{problem}%
    Знайдіть силу взаємодії двох різнойменно заряджених металевих пластин, що розташовані на
    малій відстані одна від одної. Поверхнева густина заряду пластин $\pm\sigma$, площа
    пластин $S$.

    Якщо пластини мають форму дисків радіуса $R$, чому буде дорівнювати сила взаємодії між ними на
    відстані $\ell \gg R$?
\end{problem}

%=========================================================
\begin{problem}%
    Металева сфера радіуса $R$ складена з двох півсфер. Знайдіть силу $F$, з якою відштовхуються
    ці півсфери, якщо повний заряд сфери дорівнює~$Q$.

    %Відповідь: $F= \frac{Q^2}{8R^2}$.
\end{problem}
\endgroup


\begin{Questions}
    \begin{quizlist}
        \item Що таке провідник з мікроскопічної точки зору? Чому в металах кажуть про \enquote{колективізовані} електрони і чим вони відрізняються від електронів, локалізованих на іонах?
        \item У чому полягала ідея дослідів Толмена і Стюарта? Як за напрямком відхилення стрілки гальванометра й виміряним питомим зарядом $q/m$ довели, що носіями струму в металах є саме електрони?
        \item Чому в стаціонарному стані поле всередині провідника дорівнює нулю? Як з теореми Гаусса випливає, що об'ємна густина заряду в провіднику також дорівнює нулю і всі вільні заряди розташовуються лише на поверхні?
        \item Виведіть граничні умови $E_n = 4\pi\sigma$, $E_\tau = 0$ на поверхні провідника. Що вони означають для напрямку силових ліній поля біля поверхні і~де ці лінії закінчуються?
        \item Чому поверхня провідника є еквіпотенціальною, а потенціал в усьому його об'ємі сталий?
        \item Як влаштована електрофорна машина Вімсхерста і на якому фізичному принципі (перерозподіл зарядів під дією поля, позитивний зворотний зв'язок) ґрунтується наростання заряду з кожним обертом?
        \item Чи з'являються індуковані заряди на внутрішній поверхні провідника, якщо в порожнині всередині нього немає зарядів? Доведіть це, спираючись на теорему про циркуляцію поля.
        \item Чому поле в~порожнині замкненої провідної оболонки не залежить від зовнішніх зарядів? Що змінюється, якщо в~порожнині є заряд, і~навіщо оболонку заземлюють? У чому суть електростатичного захисту (клітки Фарадея) і чому автомобіль чи літак відносно безпечні під час грози?
        \item Опишіть принцип роботи електростатичного генератора Ван де Граафа. Чому заряд, занесений стрічкою всередину сфери, повністю переходить на її зовнішню поверхню?
        \item Виведіть співвідношення $\sigma \sim 1/R$ і $E\sim1/R$ для двох з'єднаних провідним дротом куль різного радіуса за однакового потенціалу. Чому цей результат узагальнюється на залежність густини заряду від локального радіуса кривизни поверхні провідника довільної форми?
        %	\item Що таке іонний вітер і як загострений провідник у сильному полі іонізує повітря навколо себе?
        %	\item Як влаштований громовідвід Франкліна? Поясніть два механізми його захисної дії --- послаблення заряду хмари та перехоплення розряду, і чому сучасні дослідження ставлять під сумнів переважання першого з них.
        \item Сформулюйте теорему єдиності розв'язку електростатичної задачі (рівняння Пуассона чи Лапласа з граничними умовами). Які граничні умови задають на поверхнях провідників?
        \item У чому полягає метод електричних зображень і на якому факті (теоремі єдиності) він ґрунтується?
        \item Виведіть конфігурацію заряду-зображення $q'=-q$ для точкового заряду поблизу нескінченної заземленої провідної площини. Чому ця конфігурація відтворює правильне поле лише в тій половині простору, де знаходиться реальний заряд?
        \item Як метод зображень застосовується для заряду поблизу заземленої провідної сфери? Запишіть формули для величини $q'$ і положення $a$ заряду-зображення через $R$, $b$ і $q$.
        \item Чим відрізняється задача про заряд поблизу ізольованої зарядженої сфери від задачі із заземленою сферою? Скільки зарядів-зображень потрібно в цьому випадку і яка умова на їхню суму?
        \item Сформулюйте задачу про металеву кулю в однорідному зовнішньому полі $\Efield_0$. Які умови визначають розв'язок і~яку роль відіграє стрибок нормальної складової поля?
        \item Як шукають поле зовні кулі у вигляді суми однорідного поля та поля точкового диполя? Виведіть наведений дипольний момент $\vect p = R^3\Efield_0$ із граничної умови $E_\theta(R)=0$.
        \item Виведіть розподіл поверхневого заряду $\sigma(\theta) = \dfrac{3}{4\pi}E_0\cos\theta$ на металевій кулі в однорідному полі. Де поле біля поверхні кулі найбільше і~в~скільки разів воно перевищує $E_0$?
        \item Виведіть формулу тиску електричного поля на поверхню провідника $f = E^2/(8\pi)$. Чому сила дорівнює $\frac12\sigma E$, а не $\sigma E$? Куди напрямлена ця сила?
    \end{quizlist}
\end{Questions}